% UBT-AI-PROVENANCE-BEGIN
% schema: ubt-ai-provenance/v1
% tier: C_working
% ai_assistance: disclosed
% human_review: risk-based
% editorial_responsibility: Ing. David Jaroš
% policy: ../../../AI_PROVENANCE.md
% notice: Working material; exhaustive human review is not claimed.
% UBT-AI-PROVENANCE-END

\chapter{Pokračování do komplexního času a Wardovy identity}
\label{ch:ct-scheme}

\section{Proč tato kapitola zůstává v učebnici}
Starší alfa práce používala ,,CT schéma'' postavené na analytickém pokračování
\(\tau=t+i\psi\). Archivovaná implementace není aktuálním zdrojem
záznam, ale základní matematika --- analytické pokračování, měřidlo
identity a limit v reálném čase --- je standardní a užitečný.  Tato kapitola
proto učí tyto přísady a pak přesně uvádí, co dělají a dělají
nezaložit pro UBT.

\section{Komplexní čas jako parametr}
Pro analytickou veličinu $F(\tau)$,
\[
 \tau=t+i\psi
\]
umožňuje pohyb mezi oscilačními a tlumenými popisy.  Pro energii
vlastní mód, např.
\[
 e^{-iE\tau/\hbar}=e^{-iEt/\hbar}e^{+E\psi/\hbar},
\]
se znaménkem tlumení určeným konvencí pro imaginární
směr.  V theta konvenci použité později,
\[
 e^{-\pi(\psi-it)n^2}
 =e^{-\pi\psi n^2}e^{i\pi t n^2}.
\]
Nic v této analytické identitě samo o sobě nevytváří novou renormalizaci
konstantní nebo fixuje fyzickou vazbu.

\section{Wardova--Takahashiho identita: standardní argument}
V QED nechť je $S(p)$ úplným fermionovým propagátorem a
$\Gamma^\mu(p+q,p)$ plný elektromagnetický vrchol.  Invariance měřidla dává
Ward -- identita Takahashi
\[
 q_\mu\Gamma^\mu(p+q,p)
 =S^{-1}(p+q)-S^{-1}(p).
\]
Vezmeme-li limitní výnosy $q\to0$
\[
 \Gamma^\mu(p,p)=\frac{\partial S^{-1}(p)}{\partial p_\mu}.
\]
Na úrovni renormalizačních konstant se to prosazuje
\[
 \boxed{Z_1=Z_2,}
\]
za předpokladu, že regulátor a předpis pro odečítání zachovají symetrii měřidla.
Toto je výsledek standardní kalibrační teorie.  Může být použito komplexní časové pokračování
pouze po prokázání, že jeho obrys/regulátor zachovává hypotézy
vyžaduje identita.

\section{Transverzalita vlastní energie fotonu}
Z toho vyplývá stejná symetrie měřidla
\[
 q_\mu\Pi^{\mu\nu}(q)=0,
\]
takže vakuová polarizace má tvar
\[
 \Pi^{\mu\nu}(q)
 =\left(q^2g^{\mu\nu}-q^\mu q^\nu\right)\Pi(q^2)
\]
(až do konvence).  Proto může být vyjádřena renormalizace náboje
členy příčné fotonové dvoubodové funkce.  Opět se jedná o standard
omezení teorie pole, nikoli teorém specifický pro UBT.

\section{Kontrolovaný UBT most}
Bylo by nutné vytvořit komplexní renormalizační konstrukci UBT, např
jeho skutečná akce a obrys:
\begin{enumerate}
 \item Analytika dostatečná pro pokračování mezi vybraným obrysem $\psi$
       a řez v reálném čase;
 \item zachování příslušných identit měřidla/BRST;
 \item regulační a protitermické schéma kompatibilní s těmito identitami;
 \item kontinuita na standardní QED v příslušném limitu reálného času;
 \item explicitní výpočet jakéhokoli konečného CT-specifického zbytku spíše než
       nastavení podle definice.
\end{enumerate}
Teprve po těchto krocích mohl být z modelování propagován faktor specifický pro CT
výběr k odvozené veličině.

\section{Proč se archivní text \(\mathcal R_{\rm UBT}=1\) nepřebírá}
Historická větev alfa uvedla základní linii se dvěma smyčkami
\(\mathcal R_{\rm UBT}=1\) za předpokladů označených A1--A3 a následně označila
výsledná alfa základní čára ,,bez fit''.  Aktuální dokumenty o stavu Tier-A č
déle umožnit, aby toto znění stálo jako dokončená alfa derivace:
součinitelový můstek G137-B je otevřený a efektivní multiplicita má audit
mezera.  Živá učebnice proto zachovává Ward-identity matematiky, ale ano
znovu nepoužívat archivovanou větu jako aktuální zdroj záznamu.

Toto je užitečný příklad obecného pravidla: správná podmíněná teorie pole
identita může přežít, i když bude vybudována větší fenomenologická interpretace
navíc je snížená klasifikace.

\section{Vazba na theta kapitolu}
Redukované jádro theta
\[
 S(t,\psi)=\sum_n a_n e^{-\pi\psi n^2}e^{i\pi t n^2}
\]
poskytuje čistší kontrolovanou laboratoř po složitou dobu než stará alfa
renormalizační cesta.  Stejný parametr $\psi$ funguje jako proměnná doby ohřevu,
zatímco $t$ nese oscilační fázi.  Fourier, Laplace, Mellin a
Poissonovo/Jacobiho transformace pak mohou být odvozeny přesně před jakýmkoli fyzickým UBT
nárok je k nim připojen.
